\section{Descent of the actual Type I point to the original field}
\label{sec:cubic-original-field-descent}

This section supplies the original-field automorphy missing from
AC-G110. It uses the actual selected host, whose extension degrees
are prime to three. A fixed central character removes the finite
character ambiguity at each cyclic descent step. No automorphic
counting conclusion, cyclotomic level condition, or Type II
Selmer vanishing is inferred from this descent.

\paragraph{Data and retained inputs.}
Put \(p=269\), \(K=\mathbb Q(\omega)\), and let \(\chi\)
be the fixed cubic character. Use an actual Type I algebraic point
of AC-G65--AC-G66 and AC-G110, with multiplier-one restriction
\[
 \rho:G_K\longrightarrow\operatorname{GL}_3(\overline{\mathbb Q}_p),
 \qquad \nu=\rho c_p^{-1},\qquad j\in\{1,2\}.
 \tag{G111.1}
\]
Here \(c_p\) is the actual cyclotomic character, evaluated as
\(q^{-1}\) on geometric Frobenius at a prime of norm \(q\).
All the exact support and original datum hypotheses of AC-G110
are retained. In particular \(\rho\) is absolutely irreducible
on every open subgroup, and \(\nu|_{G_H}\) is attached to a
regular algebraic conjugate self-dual cuspidal representation
over the actual support host \(H\).

The chosen host has a chain of CM fields
\[
 K=H_0\subset H_1\subset\cdots\subset H_s=H,
 \qquad H_i/H_{i-1}\text{ cyclic of degree }2\text{ or }269.
 \tag{G111.2}
\]
This is the host of AC-G89/G93, followed by the selected-place
and cyclotomic steps of AC-G93-fill. The initial host is the
specified abelian extension of degree a power of two; refine
that abelian extension into quadratic steps. All subsequent
steps have degree a power of \(269\), and the specified cyclic
extensions can be refined into degree-\(269\) steps.
The cubic field used to split coefficient characters in the
class-group argument is an auxiliary field: it is not inserted
into (G111.2). In particular the explicit degree assertion in
AC-G93-fill, that the host degree has only prime factors
\(2,269\), is preserved. No descent through a degree-three
extension is used below.

We retain precisely the cyclic cuspidal descent and fiber
statements in \path{BaseChangeUnits/CycCOfCyclic.lean},
\texttt{CycC}, as supplied by the already retained
\texttt{HypothesesH'.cyc}. Its second clause descends a regular
cuspidal Satake system invariant over a cyclic prime-degree
extension; its third clause identifies any two cuspidal
descents up to a norm-trivial character of that prime order.
We also retain strong multiplicity one, global class field
theory, and the published Galois representation theorem for
regular algebraic polarized cuspidal representations over CM
fields, with its literal \texttt{AttC} normalization. These
are the same named published inputs as before. The containing
source declarations and applied wrapper proofs were read;
this is not a new proof of cyclic base change or of the
construction of those Galois representations.

\begin{lemma}[AC-G111-det: the determinant fixed on the original field]
\label{AC-G111-det}
For the actual point in (G111.1),
\[
                    \det\rho=\chi^{2j},
             \qquad \det\nu=\chi^{2j}c_p^{-3}.
 \tag{G111.3}
\]
Let \(\delta\) be the finite Hecke character of \(K\)
corresponding to \(\chi^{2j}\), using geometric reciprocity.
Then \(\delta^c=\delta^{-1}\). For every CM field
\(L\) in (G111.2), any cuspidal representation attached to
\(\nu|_{G_L}\) has central character
\(\delta_L=\delta\circ N_{L/K}\).
\end{lemma}
\begin{proof}
At the coefficient prime the actual full ordered flag has
inertia characters
\[
             \chi^j\psi c_p^e,\quad 1,\quad
                         \chi^j\psi^{-1}c_p^{-e}.
 \tag{G111.4}
\]
Their product is \(\chi^{2j}\), without a condition on the
unramified parts of the three quotient characters.

At an away prime, the modified killing condition makes the
inertia representation factor through the finite image of
the residual tuple. That image has order prime to \(p\).
Indeed the kernel of this finite residual image is exactly
among the killed inertia elements, so this is a quotient
statement, not a claim deduced only from finite inertia.
The determinant is then a character of a finite group of
order prime to \(p\). Its reduction is the product of the
residual tuple, namely \(\bar\chi^{2j}\). Roots of unity
of order prime to \(p\) have unique lifts with a prescribed
residue: if two such roots have the same residue, their
quotient lies both in the prime-to-\(p\) torsion and in
the pro-\(p\) group \(1+\mathfrak m\), hence is one.
Consequently the determinant equals \(\chi^{2j}\) on
all away inertia as well. This includes primes where the
rational character \(\psi\) is ramified.

Thus \(\det\rho\,\chi^{-2j}\) is unramified at every
finite prime of \(K\). The original point has coefficients
in one finite extension of \(\mathbb Q_p\), so this is a
continuous abelian character. Global reciprocity identifies
the maximal everywhere-unramified abelian quotient for this
imaginary quadratic field with its ideal class group.
The Eisenstein integer ring is norm-Euclidean, so
\(\operatorname{Cl}(K)=1\). The character is trivial.
The second equality in (G111.3) follows by taking the
determinant of the scalar twist.

For this cubic Kummer character, complex conjugation sends
\(\chi\) to \(\chi^{-1}\); hence the finite Hecke
character has \(\delta^c=\delta^{-1}\).
At a common unramified prime of norm \(q\), the literal
rank-three \texttt{AttC} condition says that the geometric
Frobenius eigenvalues of \(\nu\) are \(q\) times the
Satake parameters. Their product is therefore
\(q^3\omega_\Pi(v)\). On the other hand (G111.3) gives
\(q^3\delta_L(v)\). Cancellation and rank-one strong
multiplicity one give \(\omega_\Pi=\delta_L\).
This explicitly accounts for the exponent \(-3\) in
(G111.3); no reciprocal central character is substituted.
\end{proof}

\begin{theorem}[AC-G111-descent: exact conjugate self-dual descent]
\label{AC-G111-descent}
Under the stated actual hypotheses and retained published
inputs, there exists a regular algebraic conjugate self-dual
cuspidal automorphic representation \(\Pi_K\) of
\(\operatorname{GL}_3(\mathbb A_K)\) with
\[
  \operatorname{AttC}(\iota,\rho c_p^{-1},\Pi_K),
              \qquad \omega_{\Pi_K}=\delta.
 \tag{G111.5}
\]
Thus the original-field representation \(\rho c_p^{-1}\)
satisfies the literal \texttt{AutCsd0} condition.
At each step of (G111.2), its cuspidal descent with the
prescribed central character is unique up to isomorphism.
\end{theorem}
\begin{proof}
We give a complete descent step. Write \(M/L\) for one
cyclic step, of prime degree \(q\in\{2,269\}\), and
suppose \(\Pi_M\) is already regular algebraic,
cuspidal and conjugate self-dual, attached to
\(\nu|_{G_M}\), with central character \(\delta_M\).
The assertion at the top field is AC-G110-point and
AC-G111-det.

\emph{1. The actual Satake system descends.}
For all primes outside one finite set, the characteristic
polynomial of \(\nu(\operatorname{Frob}_w)\) is the
polynomial of the corresponding power of a Frobenius
element over \(L\). Primes of \(M\) over the same prime
of \(L\) have the same residue degree. Since \(\nu\)
is defined over \(G_L\), conjugation does not change
these characteristic polynomials. After dividing the
geometric eigenvalues by the common norm factor, the
Satake systems of \(\Pi_M\) at such primes agree.
This is exactly the invariance premise in the second
clause of \texttt{CycC}, not an assertion of invariance
based only on residual traces. That clause gives a regular
algebraic cuspidal descent \(\Pi_L^0\). Regularity follows
also directly from its prescribed infinity parameters:
every embedding of \(L\) extends to \(M\), where those
parameters are regular.

\emph{2. Fix the central character by a unique cubic root.}
Let \(\omega_0\) be the central character of
\(\Pi_L^0\), and put
\(\eta=\omega_0\delta_L^{-1}\).
The weak base-change identities and the product of the
Satake parameters show that
\(\eta\circ N_{M/L}=1\). For completeness, at almost
all primes \(w\) this pullback has value
\(\eta(v)^{f(w/v)}=1\); rank-one strong multiplicity
one makes the pullback character trivial globally.
By class field theory the kernel of this norm pullback
is the character group of \(\operatorname{Gal}(M/L)\).
In particular \(\eta^q=1\).

Choose an integer \(a\) with \(3a\equiv-1\pmod q\)
and set \(\mu=\eta^a\). Then
\(\mu^3=\eta^{-1}\), and \(\mu\) is norm-trivial.
Set
\[
                    \Pi_L=\Pi_L^0\otimes\mu.
 \tag{G111.6}
\]
Twisting a rank-three representation multiplies its central
character by \(\mu^3\). Hence
\(\omega_{\Pi_L}=\delta_L\). The finite-order twist
does not alter the infinity parameters, cuspidality, or
the base change to \(M\). Uniqueness of the cubic root
on this finite character group uses \(3\nmid q\).

\emph{3. The normalized descent is conjugate self-dual.}
The conjugate contragredient \(\Pi_L^{c,\vee}\) is
another cuspidal descent of \(\Pi_M\): base change
commutes with these operations, and the representation
over \(M\) is conjugate self-dual. The CM conjugations
restrict compatibly from \(M\) to \(L\).
Its central character is
\((\delta_L^c)^{-1}=\delta_L\).
The third clause of \texttt{CycC}, followed by strong
multiplicity one, says that the two descents differ by
a norm-trivial character \(\zeta\), with \(\zeta^q=1\).
Equality of their central characters gives
\(\zeta^3=1\). Thus \(\zeta=1\), and
\(\Pi_L^{c,\vee}\simeq\Pi_L\).
The same argument proves uniqueness of any cuspidal
descent with this prescribed central character. It works
also when \(q=2\); no square root in a quadratic
character group is being requested.

\emph{4. Identify the exact Galois representation.}
Apply the retained Galois representation theorem to the
regular algebraic conjugate self-dual cusp \(\Pi_L\).
Let \(r_L\) be its semisimple representation with the
literal native geometric-Frobenius normalization.
Weak base change and the attachment of \(\Pi_M\) imply
that \(r_L|_{G_M}\) and \(\nu|_{G_M}\) have the same
characteristic polynomials on almost all Frobenius
elements. Chebotarev and Brauer--Nesbitt identify their
semisimplifications. Since \(\nu|_{G_M}\) is absolutely
irreducible by AC-G110-irr, the restriction of \(r_L\)
is itself irreducible and isomorphic to it. Conjugate
the chosen matrices so that the restrictions agree.

For \(g\in G_L\), the matrix
\(r_L(g)\nu(g)^{-1}\) commutes with \(\nu(G_M)\):
this follows by conjugating any \(h\in G_M\), using
normality of \(G_M\) and equality of the two restrictions.
Schur's lemma therefore makes this matrix a scalar
\(\alpha(g)\). The representation identities give
\(\alpha(gh)=\alpha(g)\alpha(h)\), and it is one on
\(G_M\). Thus \(\alpha\) is a character of the cyclic
group of order \(q\), with \(\alpha^q=1\).

At almost all primes of \(L\), the attachment convention
and \(\omega_{\Pi_L}=\delta_L\) give
\[
           \det r_L=\delta_{L,p}c_p^{-3}
                      =\det(\nu|_{G_L}).
 \tag{G111.7}
\]
Here \(\delta_{L,p}=\chi^{2j}|_{G_L}\); equality of the
continuous characters follows from Chebotarev. Since
\(r_L=\nu|_{G_L}\otimes\alpha\), equation (G111.7)
gives \(\alpha^3=1\). Coprimality again gives
\(\alpha=1\). Therefore \(r_L\simeq\nu|_{G_L}\),
with no surviving finite twist. This proves the asserted
attachment at the next field.

Iterate the four steps down (G111.2). At the final field
they give exactly (G111.5). No automorphy lifting theorem
for a residually irreducible representation is used in
this descent: the only irreducibility needed here is the
already proved characteristic-zero irreducibility on
the actual finite extensions.
\end{proof}

\paragraph{Exact remaining interfaces.}
The conclusion is original-field automorphy of the actual
Type I point, under the already retained native support
inputs. It is not a claim that the original
\texttt{LevelAt} predicate holds. At \(269\), a character
factoring through \(K(\mu_M)\) has inertia of order
dividing \(268\cdot269^a\), whereas the relative
endpoint/middle inertia here has order three.
Consequently the original cyclotomic \texttt{KillsMu}
level wrapper cannot simply be reused. The separate
bounded local-type/counting construction must address
that exact difference. Likewise no assertion about the
remaining mixed Type II tangent or the two singular
regulator values is made here.

\paragraph{Source and geometric scope.}
The pinned source revision is
\par\noindent\texttt{7cbf7a0c32fca1badf05931879a37c3f8a2b03b4}.
The exact norm-character exponent and weak base-change
conventions are in \path{PublishedTheorems/CyclicBaseChangeC.lean},
\texttt{IsNormChar} and \texttt{WeakBC}. The descent and
fiber clauses are the complete \texttt{CycC} declaration
in \path{BaseChangeUnits/CycCOfCyclic.lean}, retained
through \path{AfDesc.lean}. The published polarized
Galois input and its short \texttt{of\_isCsd} wrapper are
in \path{GaloisCohomology2/GaloisRepOfPolarizedCuspidal.lean};
the source directory actually uses the Unicode subscript
\(2\), as recorded in the receipt. The exact
\texttt{algSat}, \texttt{AttachedSys}, and \texttt{AttC}
definitions are in \path{BlochKatoSelmerVanishes.lean}.
These source declarations do not replace the retained
published theorems by proofs. Every new determinant,
normalization, central-character, and finite-twist
argument used above has been given in full.
In the established SGA~1 finite-etale dictionary, the
fields in (G111.2) give cyclic covers and norm-trivial
characters give their finite rank-one local systems;
no new SGA~1 input is required for this interpretation.
